%! Operations with Polynomials — Unit 3, Lesson 3.1 — Student Packet Cover Sheet
\documentclass[10pt]{article}
\usepackage{algebra2-article}
\usepackage{algebra2-boxes}
\usepackage{ltablex}
\keepXColumns

\begin{document}

% Full-bleed forest banner
\begin{tikzpicture}[remember picture, overlay]
  \fill[forest] (current page.north west)
    rectangle ([yshift=-0.9in]current page.north east);
\end{tikzpicture}
\vspace{-0.8in}

\begin{center}
  {\color{white}\textbf{\LARGE Algebra 2}}\\[4pt]
  {\color{forestmid}\large\textbf{Unit 3 \quad Polynomial Functions}}\\[6pt]
  {\color{white}\textbf{\large Lesson 3.1 \quad Operations with Polynomials}}
\end{center}

\vspace{0.22in}
\namedateperiod
\vspace{0.18in}

\begin{learningtargetbox}
\small
\begin{enumerate}[leftmargin=1.5em, itemsep=5pt]
  \item \ldots{} write a polynomial in \textbf{standard form} and name its \textbf{degree}, \textbf{leading coefficient}, and number of \textbf{terms}.
  \item \ldots{} \textbf{add} and \textbf{subtract} polynomials by combining like terms, distributing a subtraction sign to \emph{every} term.
  \item \ldots{} \textbf{multiply} polynomials in one and two variables --- monomial $\times$ polynomial, binomial $\times$ binomial, and binomial $\times$ trinomial --- and apply the \textbf{special products} $(a\pm b)^2$ and $(a+b)(a-b)$.
  \item \ldots{} predict a product's \textbf{degree} and \textbf{leading coefficient} and use them to state its \textbf{end behavior}.
\end{enumerate}
\end{learningtargetbox}

\vspace{0.14in}

\begin{tocbox}
\small
\renewcommand{\arraystretch}{1.5}
\begin{tabularx}{\linewidth}{@{} c l X r @{}}
\rowcolor{forestbg}
\textbf{\#} & \textbf{Component} & \textbf{Description} & \textbf{Score} \\
\midrule
1 & Warm-Up        & Spiral review: powers, like terms, distributing, and FOIL                     & \blank{1.2cm} \\
2 & Guided Notes   & Adding, subtracting, and multiplying polynomials; special products             & \blank{1.2cm} \\
3 & Group Activity & \emph{Building Polynomials} --- products, predictions, and models, by tier     & \blank{1.2cm} \\
4 & Exit Ticket    & Quick check: a subtraction, a product, and an SOL-style item                   & \blank{1.2cm} \\
5 & Homework       & Mixed operations, special products, and an area-model extension                & \blank{1.2cm} \\
\midrule
  & \multicolumn{2}{r}{\textbf{Total}} & \blank{1.2cm} \\
\end{tabularx}
\end{tocbox}

\vspace{0.10in}

\begin{remindbox}
\small Every polynomial operation comes from \textbf{one} idea: \textbf{like terms combine, unlike
terms do not}. \textbf{Adding} and \textbf{subtracting} only \emph{collect} terms --- and a minus
sign in front of a parenthesis is a factor of $-1$, so it flips the sign of \emph{every} term inside.
\textbf{Multiplying} \emph{creates} terms: multiply \emph{every} term by \emph{every} term
(coefficients multiply, exponents \emph{add}), then collect. Three products are worth knowing on
sight: $(a+b)^2=a^2+2ab+b^2$, $(a-b)^2=a^2-2ab+b^2$, and $(a+b)(a-b)=a^2-b^2$ --- and note that
$(a+b)^2$ is \emph{never} $a^2+b^2$. Finally, a product's shape is predictable: \textbf{degrees add
and leading coefficients multiply}, so you know a product's end behavior before you expand it.
\end{remindbox}

\end{document}
